\section{What is Musical Notation?}
Musical notation can be thought of as a graph, where the \emph{X axis}, the horizontal axis, is musical time that progresses from the first note to the last, and the \emph{Y axis}, the vertical axis, represents musical pitch, which tells which notes to play. The musical alphabet is composed of seven letter names from \mnote{A} through \mnote{G} which are placed on either a line or a space of the \emph{musical staff}. The musical staff itself, which is a collection of five horizontal lines, doesn't tell us which note to play, a \emph{clef} is required to tell us where \emph{middle \mnote{C}} is, a where every other note is by extension. For the purposes of this book, and the vast majority of musical applications, only the \emph{treble clef}: \clefG{} and \emph{bass clef}: \clefF{} are required. The following examples show the notes of each of the clefs with labels, for treble clef, the middle \mnote{C} is below the staff, and for bass clef the middle \mnote{C} is above the staff:

\bookExample{Notes of the Treble Clef}{./excerpts/treble-clef.ly}

\bookExample{Notes of the Bass Clef}{./excerpts/bass-clef.ly}

The X axis of music tells us how long to play each note. For the most part, our examples will be relatively simple from a rhythmic perspective. \emph{Whole notes}: \wholeNote{} last for four beats, \emph{half notes}: \halfNote{} last for two beats, \emph{quarter notes}: \quarterNote{} last for one beat, \emph{8th notes} \eighthNote{} last for half a beat, and \emph{16th notes}: \sixteenthNote{} last for a quarter of a beat. Notice that eighth notes, and every rhythmic value faster than that has a flag. They can be beamed together if they are in the same beat. We will see many examples where 8th notes are beamed together like so: \twoBeamedQuavers{}. The convention we use for this book is to beam in half notes. We will often see 4 8th notes beamed together, or 8 sixteenth notes, and so on.

But musical notes aren't sounding all of the time, we need a way to indicate silence. For this we use symbols called rests. They have rhythmic values that correspond to the above notes. We have the \emph{whole rest}: \wholeRest{}, the \emph{half rest}: \halfRest{}, the \emph{quarter rest}: \quarterRest{}, the \emph{eighth rest}: \eighthRest{}, and the \emph{sixteenth rest}: \sixteenthRest{}.

Without the concept of meter in music, we would just see a very large sequence of notes one after the other, without deliniation. The meter that our music is in tells us how many notes, and what kind of notes, fit inside the bar. For the vast majority of jazz tunes, the meter will be \meter{4}{4}. We can see it in context here, to the right of the treble clef. So, in \meter{4}{4} meter, we have four quarter notes per bar.

Our musical system has the ability to alter the pitch of a given note with special symbols called \emph{accidentals} where the \emph{sharp symbol}: \sharp{}, raises the note by one \emph{half step} and the \emph{flat symbol}: \flat{}, lowers the pitch of the note by one half step and the \emph{natural symbol} \natural{}, cancels the effect of the existing accidental. It's important to note that the accidental only has its effect for the measure it's inside of and doesn't cross the bar line. This means that we have to notate an accidental everytime we need it.

Key signatures are a default set of accidentals that are applied to each bar of a piece where a given key signature is active. Here we will see the key signature that corresponds to the key \mnote{F}, it has a single flat applied to the note \mnote{B} Which means that for every bar where the active key signature is \mnote {F}, every \mnote{B} becomes a \mnote[flat]{B} unless an accidental overrides it. Any accidental symbol is only active for a single bar, then the key signature returns to its normal behavior:

\bookExample{Key Signature Example}{./excerpts/key-example.ly}
